\section{Resultados e Discuss\~ao}
\label{sec:resultados}

\subsection{Evidência das UDVs}
\label{sec:resultados-udv}

Os controles disponíveis indicam fragilidade no vínculo por similaridade semântica: nas 183 opiniões
de treino com citação confiável, a sentença de maior cosseno com a opinião completa é a da própria
citação em 112 (61,2\%). O intervalo de 95\% dessa proporção, por \textit{bootstrap} de audiências,
vai de 53,0\% a 68,8\%. Com os trechos entre aspas removidos da opinião, a sentença da citação é a
de maior cosseno em 12 de 142 consultas do treino e em 4 de 40 da validação. Esses números mostram
que a sentença escolhida por similaridade difere com frequência da passagem citada, mas não separam
o erro de outra sentença que também sustente a opinião. Essa separação cabe à validação humana, na
qual <resultado da validação humana: precisão estrita da citação literal e tolerante do nível
\texttt{semantic\_match\_high}, com intervalos de Wilson, critérios atingidos ou não, e concordância
do anotador consigo mesmo>.

\subsection{Simulação}
\label{sec:resultados-simulacao}

\input{tables/main-results}

Na validação, <n> das <n> opiniões ligadas geraram perguntas, de <n> atores em <n> audiências. O
maior acerto da condição 2 foi obtido com $k = {<}k{>}$, e o acerto de cada valor da grade é
reportado em <local>. No teste, <n> das 106 opiniões ligadas
geraram perguntas, de <n> atores em <n> audiências (Tabela~\ref{tab:resultados}). A soma das
probabilidades das quatro letras teve média <0,00>, <0,00> e <0,00> e mínimo <0,00>, <0,00> e
<0,00> nas condições 0, 1 e 2.

<Discussão dos resultados da simulação, a escrever com os números da execução.>

% Pendente: preencher a Figura fig:simulacao com uma pergunta de teste (itens do perfil, trechos recuperados, as quatro opções e as probabilidades nas condições 0 a 2), de um ator que não tenha entrado no teste preliminar; a regra de escolha da pergunta ainda precisa ser fixada antes da rodada final.

\subsection{Limitações}
\label{sec:limitacoes}

\textbf{Gabarito.} A opção correta é a opinião que a matéria atribui à pessoa, na redação da
matéria, e o acerto mede o reconhecimento dessa atribuição, que é uma seleção editorial do que foi
dito. Os distratores podem coincidir com posições que a pessoa também sustenta, sobretudo nas
condições que informam o lado dela, e essas perguntas admitem mais de uma resposta plausível. A
frequência dessas perguntas não foi medida.

\textbf{Informação comum às condições.} O nome e o cargo entram em todas as condições, e o cargo às
vezes registra a função da pessoa na própria audiência avaliada, como relator ou presidente da
comissão. O modelo pode associar ao nome posições conhecidas fora do conjunto. A condição 0 contém
esses efeitos e os do assunto e da redação das opções. As diferenças entre condições os cancelam
somente se eles se somarem ao efeito do material acrescentado, sem interagir com ele.

\textbf{Pré-treinamento.} A data de corte do pré-treinamento do modelo simulador, <data de corte>, é
<anterior ou posterior> ao período de teste. Se for posterior, as matérias das audiências de teste
podem ter sido vistas no pré-treinamento, com efeito em todas as condições.

\textbf{Tamanho da amostra.} As opiniões de teste ligadas a atores com perfil são no máximo 106, em
27 audiências, e um intervalo que contém o zero indica que a diferença não foi detectada com essa
amostra. O valor de $k$ é escolhido numa amostra de validação pequena e agrupada por audiência, em
que diferenças de poucas perguntas entre valores da grade podem refletir ruído.

\textbf{Desenho da avaliação.} Como as condições são encadeadas, a recuperação sem o perfil não é
avaliada, e o ganho do perfil sobre a simples recuperação das falas não é medido. A avaliação
também não separa assuntos já tratados nas falas de treino de assuntos novos.

\textbf{Perfis.} Como o modelo simulador é o mesmo LLM que escreveu os perfis, o efeito do perfil é
medido com um leitor que compartilha o vocabulário e os vieses de quem o redigiu. Esse efeito pode
não se repetir com outro modelo simulador. A fidelidade dos perfis às falas não foi medida por
anotação humana.

\textbf{Evidência da UDV.} A precisão da associação entre o nome citado na matéria e os turnos da
pessoa não foi medida. O corte de 0,45 entre \texttt{semantic\_match\_high} e
\texttt{semantic\_match\_weak} não foi validado como medida de confiança, e a variante da calibração
com negativos da própria pessoa daria 0,50 (Apêndice~\ref{ap:calibracao}). Nem a construção nem o
verificador avaliam se a evidência sustenta a opinião, o que cabe à validação humana
(Seção~\ref{sec:setup-udv}).
